\documentclass[english, parskip=full]{scrartcl}
\usepackage[utf8]{inputenc}  % Kodierung der Datei
\usepackage[english]{babel}  % Sprache auf Deutsch 
\usepackage{color}
\usepackage{graphicx, gensymb}
\usepackage{amsfonts, cancel, amssymb}
\usepackage{amsmath, amsthm, array, esint}
\usepackage{booktabs, multirow}  % schönere Tabellen
\usepackage{caption}  % Anpassung der Captions
\usepackage{scrlayer-scrpage}    % Manipulation des Seitenstils
\pagestyle{scrheadings}  % Stil für die Seite setzen
\automark{section}  % Abschnittsnamen als Seitenbeschriftung 
\setheadsepline{.5pt}    % Kopzeile durch Linie abtrennen
\usepackage[hidelinks]{hyperref}  % Links und weitere PDF-Features
\usepackage{typearea, tikz, pgffor, verbatim}
\usetikzlibrary{calc, arrows.meta,arrows, graphs, quotes, positioning}
\usepackage[cache=false, outputdir=build]{minted}
\usemintedstyle{vs}
\usepackage{placeins} % provides \FloatBarrier

\definecolor{bg}{rgb}{0.9,0.9,0.9}
\newcommand\numberthis{\addtocounter{equation}{1}\tag{\theequation}}
\usepackage{svg}
\usepackage{siunitx}
\usepackage{physics}
\usepackage{tensor}
\renewcommand{\bra}{}
\newcommand{\kom}{}
\newcommand{\brak}{}
\renewcommand{\braket}{}
\newcommand{\intx}{}
\newcommand{\intr}{}
\newcommand{\kla}{}
\newcommand{\klam}{}
\newcommand{\klammer}{}
\newcommand{\oper}{}
\newcommand{\tecxt}{}
\newcommand{\Bra}{}
\newcommand{\Ket}{}
\renewcommand{\i}{\text{i}}
\newcommand{\e}{\text{e}}
\newcommand*{\Fourier}{\textsc{Fourier}\ }
\newcommand*{\F}{\mathcal{F}}
\newcommand*{\sinc}{\text{sinc\,}}
\newcommand{\conv}{\mathop{\scalebox{3.5}{\raisebox{-0.38ex}{\makebox[0.5\width][c]{$\ast$}}}}\limits}

\newcommand*{\titel}{Special solutions to the Schrödinger equation}
\newcommand*{\autor}{Joachim Schwardt}

\hypersetup{pdfauthor={\autor}, pdftitle={\titel}}

%\titlehead{titlehead}
\subject{Quantum Theory I}
\title{\titel}
%\author{\autor}
\date{\today}

\begin{document}

%\begin{titlepage}
\maketitle  % Titel setzen
%\tableofcontents  % Inhaltsverzeichnis setzen
%\end{titlepage}

\section{One dimensional problems}
\subsection{Harmonic oscillator}
The harmonic oscillator potential is $V(x) = \frac{m\omega^2}{2}x^2$.
This is the most famous potential, and is usually solved in an algebraic fashion using creation and annihilation operators.
Here we will solve it using the standard methods of ordinary differential equations used throughout this chapter.
The problem at hand is
\begin{align}
E\psi(x) &= -\frac{\hbar^2}{2m}\psi''(x) + \frac{m\omega^2}{2}x^2\psi(x).
\end{align}
We will factor out the asymptotic behaviour using $\psi(x) := \phi(x) \e^{-\gamma x^b}$ for some constants $\gamma,b>0$.
The transformation of the derivative reads
\begin{align}
%\frac{\psi'}{\psi} &= \frac{\phi'}{\phi} - \gamma bx^{b-1},\\
\frac{\psi''}{\psi} &= \frac{\phi''}{\phi} - \frac{2\gamma bx^{b-1}\phi'}{\phi} + \gamma^2b^2x^{2b-2} - \gamma b(b-1)x^{b-2}.
\end{align}
Inserting into the differential equation yields
\begin{align}
0 &= \frac{2mE}{\hbar^2} - \frac{m^2\omega^2}{\hbar^2}x^2 +  \gamma^2b^2x^{2b-2} - \gamma b(b-1)x^{b-2} - \frac{2\gamma bx^{b-1}\phi'}{\phi} + \frac{\phi''}{\phi}.
\end{align}
Choosing $b=2$ and $\gamma = \frac{m\omega}{2\hbar}$ allows us to eliminate the dominant $x^2$-term.
The equation then simplifies to
\begin{align}
0 &= \frac{2mE}{\hbar^2} - 2\gamma - 4\gamma x\frac{\phi'}{\phi} + \frac{\phi''}{\phi}.
\end{align}
%Before attempting a Taylor expansion, we first eliminate a potential fractional term using the substitution $\phi(x) = x^\alpha P(x)$.
%The derivatives become
%\begin{align}
%\frac{\phi'}{\phi} &= \frac{P'}{P} + \frac{\alpha}{x},\\
%\frac{\phi''}{\phi} &= \frac{P''}{P} + \frac{2\alpha P'}{xP} + \frac{\alpha(\alpha - 1)}{x^2}
%\end{align}
%leading to the new equation
%\begin{align}
%0 &= \frac{2mE}{\hbar^2} - 2\gamma - 4\gamma x \kla\left( \frac{P'}{P} + \frac{\alpha}{x} \right) + \frac{P''}{P} + \frac{2\alpha P'}{xP} + \frac{\alpha(\alpha - 1)}{x^2}.
%\end{align}
%Choosing $\alpha = \frac{mE}{2\hbar^2\gamma} - \frac{1}{2}$ removes the energy from the equation.
%Collecting terms then yields
%\begin{align}
%0 &= \frac{\alpha(\alpha - 1)}{x^2} + \kla\left( \frac{2\alpha}{x} - 4ax \right) \frac{P'}{P} + \frac{P''}{P}.
%\end{align}
%We expand the function in a series $P(x) = \sum a_nx^n$
%\begin{align}
%0 &= \sum_{n\ge 0} \klammer\left\{ x^{n-2}a_n \klam\left[ \alpha(\alpha - 1) + 2\alpha n + n(n-1) \right] - 4\gamma a_n n x^n \right\} = \\
%&= \sum_{n\ge 0} x^n \klammer\left\{ a_{n+2} \klam\left[ \alpha(\alpha - 1) + 2\alpha (n+2) + (n+2)(n+1) \right] - 4\gamma a_nn \right\}.
%\end{align}
%Here we assumed $a_1 = 0$, since we do not allow terms of the form $\frac{1}{x}$ in the expansion.
%This gives a recurrence relation 
%\begin{align}
%\frac{a_{n+2}}{a_n} &= \frac{4\gamma n}{\alpha (\alpha - 1) + 2\alpha (n+2) + (n+2)(n+1)} \sim \frac{4\gamma}{n}.
%\end{align}
%Since $a_{2n+1} = 0$, the function $P(x)$ is symmetric and has the asymptotic form $P\sim \sum_{n\ge 0} \frac{(4\gamma)^n}{n!} x^{2n} \sim \e^{4\gamma x^2}$.
%This solution is non-normalizable, and thus the recurrence must terminate at some $N\in\mathbb{N}$.
We expand the function in a series $\phi(x) = \sum a_nx^n$
\begin{align}
0 &= \sum_{n\ge 0} \klammer\left\{ x^{n-2}a_n n(n-1) + x^n a_n \klam\left[ \frac{2mE}{\hbar^2} - 2\gamma - 4\gamma n \right] \right\} = \\
&= \sum_{n\ge 0} x^n \klammer\left\{ a_{n+2} (n+2)(n+1) - 4\gamma a_n \klam\left[ n + \frac{1}{2} - \frac{mE}{2\hbar^2\gamma} \right] \right\}.
\end{align}
There exist two solutions, which are essentially determined by $a_0=0$ and $a_1=1$ or $a_0=1$ and $a_1=0$.
In both cases we find the recurrence relation 
\begin{align}
\frac{a_{n+2}}{a_n} &= \frac{4\gamma}{n+2} \frac{ n + \frac{1}{2} - \frac{mE}{2\hbar^2\gamma}}{n+1} \sim \frac{4\gamma}{n+2}.
\end{align}
In both cases the function $P(x)$ ends up with the asymptotic form $P\sim \sum_{n\ge 0} \frac{(4\gamma)^n}{n!} x^{2n} \sim \e^{4\gamma x^2}$.
The resulting solution is non-normalizable, and thus the recurrence must terminate at some $N\in\mathbb{N}$ with $a_N = 0$.
We find a discretized number of solutions with energies
\begin{align}
E_n &= \frac{2\hbar^2\gamma}{m} \kla\left( n + \frac{1}{2} \right) = \hbar\omega \kla\left( n + \frac{1}{2} \right).
\end{align}
%The recurrence relation can then be rewritten as
%\begin{align}
%\frac{a_n^{(N)}}{a_{n-2}^{(N)}} &= \frac{2m\omega}{n\hbar} \frac{N - n}{n-1}.
%\end{align}
%Distinguishing between even and odd terms, we may rewrite this as
%\begin{align}
%\frac{a_{2n}^{(2N)}}{a_{2(n-1)}^{(2N)}} &= \frac{2m\omega}{\hbar} \frac{2N - 2n}{2n (2n - 1)}, \\
%\frac{a_{2n+1}^{(2N+1)}}{a_{2(n-1) + 1}^{(2N+1)}} &= \frac{2m\omega}{\hbar} \frac{2N - 2n}{2n (2n + 1)}
%\end{align}
%%%%For $a_0 = 1$ the first few terms are
%%%%\begin{align}
%%%%a_2^{(2N)} &= \frac{1}{2} \frac{2m\omega}{\hbar} (2N-2),\\
%%%%a_4^{(2N)} &= \frac{1}{2\cdot 4} \kla\left( \frac{2m\omega}{\hbar} \right)^2 (2N-2)(2N-4)
%%%%\end{align}
%%%%and for $a_1 = 1$ they are
%%%%\begin{align}
%%%%a_3^{(2N+1)} &= \frac{1}{3}\frac{2m\omega}{\hbar} (2N-2),\\
%%%%a_5^{(2N+1)} &= \frac{1}{3\cdot 5} \kla\left( \frac{2m\omega}{\hbar} \right)^2 (2N-2)(2N-4).
%%%%\end{align}
%%%%We may summarize these as
%%%%\begin{align}
%%%%a_{2k}^{(2N)} &= \frac{(2N-2)\dots (2N - 2k)}{(2k)!!} \kla\left( \frac{2m\omega}{\hbar} \right)^k = \frac{(2N-2)\dots (2N - 2k)}{(2k)!!} \kla\left( \frac{2m\omega}{\hbar} \right)^k, \\
%%%%a_{2k+1}^{(2N+1)} &= \frac{(N-2)\dots (N-k)}{k!!} \kla\left( \frac{2m\omega}{\hbar} \right)^{\lfloor \frac{k}{2} \rfloor}.
%%%%\end{align}
%Except for normalization, the respective eigenfunctions are
%\begin{align}
%\psi_{2n}(x) &= \e^{-\frac{m\omega}{2\hbar}x^2} \sum_{k = 0}^n x^{2k} \klam\left[ \kla\left( \frac{4m\omega}{\hbar} \right)^k \frac{(n-k)\dots (n-1)}{(2k)!} \right] = \\
%&= \e^{-\frac{m\omega}{2\hbar}x^2} \sum_{k = 0}^n \frac{(n-1)!}{(n-k-1)! (2k)!} \kla\left( \frac{4m\omega}{\hbar} x^2 \right)^k.
%\end{align}

\subsection{Morse potential}
The Morse potential is defined as $V(r) = V_0\kla\left( 1 - \e^{-a(r-r_0)} \right)^2 - V_0$, where $r_0$ is the equilibrium distance and $V(r\rightarrow\infty) \rightarrow 0$.
The Schrödinger equation reads
\begin{align}
E\psi(r) &= -\frac{\hbar^2}{2m}\psi''(r) + V(r)\psi(r).
\end{align}
It is useful to substitute $z = \e^{-a(r-r_0)}$.
The second derivative then transforms as
\begin{align}
\frac{\tecxt\text{d}^2\psi(z)}{\tecxt\text{d}r^2} &= \frac{\tecxt\text{d}^2\psi(z)}{\tecxt\text{d}z^2} \kla\left( \frac{\tecxt\text{d}z}{\tecxt\text{d}r} \right)^2 + \frac{\tecxt\text{d} \psi(z)}{\tecxt\text{d}z} \frac{\tecxt\text{d}^2z}{\tecxt\text{d}r^2} = \psi''(z)a^2z^2 + \psi'(z)a^2z
\end{align}
which results in the differential equation
\begin{align}
0 &= \frac{2m}{\hbar^2a^2} \kla\left( E - V_0z^2 + 2V_0z \right)\psi(z) + z^2\psi''(z) + z\psi'(z).
\end{align}
Define $\epsilon := \frac{E}{V_0}$ and let $\psi(z) = \phi(z) \e^{-\gamma z}$ for some $\gamma > 0$.
Then we find
\begin{align}
\frac{\psi'}{\psi} &= \frac{\phi'}{\phi} - \gamma,\\
\frac{\psi''}{\psi} &= \frac{\phi''}{\phi} + \gamma^2 - 2\gamma \frac{\phi'}{\phi}
\end{align}
and the differential equation
\begin{align}
0 &= \frac{2mV_0}{\hbar^2a^2} \kla\left( \epsilon - z^2 + 2z \right) + z^2 \kla\left( \gamma^2 - 2\gamma \frac{\phi'}{\phi} + \frac{\phi''}{\phi} \right) + z\kla\left( \frac{\phi'}{\phi} - \gamma \right).
\end{align}
Choosing $\gamma = \sqrt{\frac{2mV_0}{\hbar^2a^2}}$ eliminates the term that would be of order $z^{n+2}$ in a potential Taylor expansion of $\phi(z)$.
Collecting terms now gives
\begin{align}
0 &= \gamma^2 \kla\left( \epsilon + 2z - \frac{z}{\gamma} \right) + \frac{z\phi'}{\phi} \kla\left( 1 - 2\gamma z \right) + \frac{z^2\phi''}{\phi}.
\end{align}
Before the expansion, we will split off a polynomial factor using the transformation $\phi(z) = z^\alpha P(z)$.
This gives
\begin{align}
\frac{\phi'}{\phi} &= \frac{\alpha}{z} + \frac{P'}{P},\\
\frac{\phi''}{\phi} &= \frac{\alpha(\alpha - 1)}{z^2} + \frac{P''}{P} + \frac{2\alpha P'}{Pz}
\end{align}
and thus transforms the differential equation to
\begin{align}
0 &= \gamma^2\epsilon + 2z\gamma \kla\left( \gamma - \frac{1}{2} \right) + \kla\left( \alpha + \frac{zP'}{P} \right) \kla\left( 1 - 2\gamma z \right) + \alpha(\alpha - 1) + \frac{z^2P''}{P} + \frac{2\alpha zP'}{P} .
\end{align}
We choose $\alpha = \gamma \sqrt{-\epsilon}$ to eliminate the energy from the equation. 
Note that we consider bound states, and thus $\epsilon < 0$.
Collecting terms once again gives
\begin{align}
0 &= 2z\gamma \kla\left( \gamma - \frac{1}{2} - \alpha \right) P + \kla\left( 1 - 2\gamma z + 2\alpha \right) zP' + z^2P''.
\end{align}
At this point we may attempt to find a polynomial solution for $P(z)$, i.e. $P(z) = \sum_{n\ge 0} a_nz^n$.
Inserting gives
\begin{align}
0 &= \sum_{n\ge 0} \klammer\left\{ z^{n+1} a_n\klam\left[ -2\gamma n + 2\gamma \kla\left( \gamma - \frac{1}{2} - \alpha \right) \right] + z^na_n \klam\left[ n(n-1) + n(1+2\alpha) \right] \right\} = \\
&= \sum_{n\ge 0} z^{n+1} \klammer\left\{ -2\gamma a_n \klam\left[ n - \gamma + \frac{1}{2} + \alpha \right] + (n+1)a_{n+1} \klam\left[ n + 1 + 2\alpha \right] \right\}
\end{align}
where we performed an index shift $n\rightarrow n+1$ in the second part of the summation.
This gives the recurrence relation
\begin{align}
\frac{a_{n+1}}{a_n} &= \frac{2\gamma}{n+1} \frac{n - \gamma + \frac{1}{2} + \alpha}{n+1 + 2\alpha} \kom\overset{\substack{\text{as }{n\rightarrow\infty} \\ |}}{\sim} \frac{2\gamma}{n+1}.
\end{align}
Due to the asymptotic behaviour, we find $a_{n+1} \sim \frac{(2\gamma)^{n+1}}{(n+1)!}$, or equivalently $P(z) \sim \e^{2\gamma z}$.
This results in a non-normalizable solution, and thus the recursion must terminate at some step $N\in\mathbb{N}$
\begin{align}
0 &\overset{!}{=} N - \gamma + \frac{1}{2} + \alpha.
\end{align}
This gives an equation for the energy levels $\epsilon_N$ of the form
\begin{align}
\epsilon_N &= -\frac{\alpha^2}{\gamma^2} = \frac{-1}{\gamma^2} \kla\left( \gamma - N - \frac{1}{2} \right)^2.
\end{align}
Reinserting the constants we find
\begin{align}
E_n &= -\frac{\hbar^2a^2}{2m} \kla\left( \sqrt{\frac{2mV_0}{\hbar^2a^2}} - n - \frac{1}{2} \right)^2
\end{align}
for $n=1,\dots, \lfloor \gamma \rfloor$.

\subsection{Anharmonic oscillator}
Here we consider the anharmonic oscillator potential $V(x) = \frac{m\omega^2}{2}x^2 + Ax^4$.
The Schrödinger equation is
\begin{align}
E\psi(x) &= -\frac{\hbar^2}{2m}\psi''(x) + \frac{m\omega^2}{2}x^2\psi(x) + Ax^4\psi(x).
\end{align}
We will factor out the asymptotic behaviour using $\psi(x) = \phi(x) \e^{-\gamma x^b}$ for some constants $\gamma,b>0$.
The transformation of the derivative reads
\begin{align}
%\frac{\psi'}{\psi} &= \frac{\phi'}{\phi} - \gamma bx^{b-1},\\
\frac{\psi''}{\psi} &= \frac{\phi''}{\phi} - \frac{2\gamma bx^{b-1}\phi'}{\phi} + \gamma^2b^2x^{2b-2} - \gamma b(b-1)x^{b-2}.
\end{align}
Inserting into the differential equation yields
\begin{align}
0 &= \frac{2mE}{\hbar^2} - \frac{m^2\omega^2}{\hbar^2}x^2 - \frac{2mA}{\hbar^2}x^4 + \gamma^2b^2x^{2b-2} - \gamma b(b-1)x^{b-2} - \frac{2\gamma bx^{b-1}\phi'}{\phi} + \frac{\phi''}{\phi}.
\end{align}
Choosing $b=3$ and $\gamma = \sqrt{\frac{2mA}{9\hbar^2}}$ allows us to eliminate the dominant $x^4$-term.
The equation then simplifies to
\begin{align}
0 &= \frac{2mE}{\hbar^2} - \frac{m^2\omega^2}{\hbar^2}x^2 - 6\gamma x - 6\gamma x^2\frac{\phi'}{\phi} + \frac{\phi''}{\phi}.
\end{align}
%%%%AGAIN:
%%%%\begin{align}
%%%%0 &= \frac{2mE}{\hbar^2} - \frac{m^2\omega^2}{\hbar^2}x^2 - 6\gamma x - 6\gamma x^2 \kla\left( \frac{\phi'}{\phi} - \kappa bx^{b-1} \right) + \frac{\phi''}{\phi} - \frac{2\kappa bx^{b-1}\phi'}{\phi} + \kappa^2b^2x^{2b-2} - \kappa b(b-1)x^{b-2}
%%%%\end{align}
%%Before attempting a Taylor expansion, we first eliminate a potential fractional term using the substitution $\phi(x) = x^\alpha P(x)$.
%%The derivatives become
%%\begin{align}
%%\frac{\phi'}{\phi} &= \frac{P'}{P} + \frac{\alpha}{x},\\
%%\frac{\phi''}{\phi} &= \frac{P''}{P} + \frac{2\alpha P'}{xP} + \frac{\alpha(\alpha - 1)}{x^2}
%%\end{align}
%%leading to the new equation
%%\begin{align}
%%0 &= \frac{2mE}{\hbar^2} - \frac{m^2\omega^2}{\hbar^2}x^2 - 6\gamma x - 6\gamma x^2 \kla\left( \frac{P'}{P} + \frac{\alpha}{x} \right) + \frac{P''}{P} + \frac{2\alpha P'}{xP} + \frac{\alpha(\alpha - 1)}{x^2}.
%%\end{align}
Let $P(x) = \sum a_nx^n$.
Then we find
\begin{align}
0 &= \sum_{n\ge 0} \klammer\left\{ x^{n-2} n(n-1)a_n + x^n \frac{2mE}{\hbar^2}a_n - 6\gamma x^{n+1}a_n (n+1) - x^{n+2} \frac{m^2\omega^2}{\hbar^2} a_n \right\} = \\
&= \sum_{n\ge 0} x^n \klammer\left\{ a_{n+2} (n+2)(n+1) + a_n \frac{2mE}{\hbar^2} -6\gamma n a_{n-1} - \frac{m^2\omega^2}{\hbar^2}a_{n-2} \right\}.
\end{align}
This yields the recurrence relation
\begin{align}
a_{n+2} &= \frac{-\frac{2mE}{\hbar^2} a_n + 6\gamma n a_{n-1} + \frac{m^2\omega^2}{\hbar^2} a_{n-2}}{(n+2)(n+1)}.
\end{align}
Consider a more general form of this relation
\begin{align}
a_{n} &= \beta_na_{n-2} + \gamma_na_{n-3} + \delta_na_{n-4}.
\end{align}
Let $b_n := a_{n-1}$, $c_n := a_{n-2}$ and $d_n := a_{n-3}$.
Then we may write this as
\begin{align}
A_n &:= \begin{pmatrix}
a_n \\ b_n \\ c_n \\ d_n
\end{pmatrix} = \begin{pmatrix}
0 & \beta_n & \gamma_n & \delta_n \\
1 & 0 & 0 & 0 \\
0 & 1 & 0 & 0 \\
0 & 0 & 1 & 0
\end{pmatrix} \begin{pmatrix}
a_{n-1} \\ b_{n-1} \\ c_{n-1} \\ d_{n-1}
\end{pmatrix} =: T_nA_{n-1}.
\end{align}



%\begin{thebibliography}{9}
%\bibitem{example} description, \url{https://www.exmapleurl}, day.\,month~2021.
%\end{thebibliography}

\end{document}